% Part: normal-modal-logic
% Chapter: frame-correspondence
% Section: introduction

\documentclass[../../../include/open-logic-section]{subfiles}

\begin{document}

\olfileid[pt]{nml}{frd}{int}

\olsection{Introdução}

Uma questão de interesse para os lógicos modais é a relação entre a relação
de acessibilidade e a verdade de certas !!{formula}s em modelos que possuem
essa relação de acessibilidade. Por exemplo, suponha que a relação de
acessibilidade seja reflexiva, O.s.\@, que, para todo $w \in W$, $Rww$. Em
outras palavras, cada mundo é acessível a partir de si mesmo. Isso significa
que, quando $\Box !A$ é verdadeira em um mundo~$w$, o próprio $w$ está entre
os mundos acessíveis nos quais~$!A$ deve, portanto, ser verdadeira. Assim, se
a relação de acessibilidade~$R$ de~$\mModel{M}$ é reflexiva, então, quaisquer
que sejam o mundo $w$ e a fórmula $!A$ que tomemos, $\Box !A \lif !A$ será
verdadeira nesse mundo (em outras palavras, o esquema $\Box p \lif p$ e todas
as suas instâncias de substituição são verdadeiros em~$\mModel{M}$).

A recíproca, contudo, é falsa. Não ocorre, p.ex., que, se
$\Box p \lif p$ é verdadeira em~$\mModel{M}$, então $R$ é reflexiva. De fato,
podemos encontrar facilmente um modelo não reflexivo~$\mModel{M}$ no qual
$\Box p \lif p$ é verdadeira em todos os mundos: tome o modelo com um único
mundo~$w$, que não é acessível a partir de si mesmo, mas no qual $w \in V(p)$.
Escolhendo adequadamente o valor de verdade de $p$, podemos tornar
$\Box !A \lif !A$ verdadeira em um modelo que não é reflexivo.

A solução é retirar da equação a atribuição de variáveis~$V$. Se exigirmos
que $\Box p \lif p$ seja verdadeira em todos os mundos de~$\mModel{M}$,
independentemente de quais mundos pertençam a~$V(p)$, então é necessário que
$R$ seja reflexiva. Pois, em todo modelo não reflexivo, haverá ao menos um
mundo~$w$ tal que não $Rww$. Se definirmos $V(p) = W \setminus \{w\}$, então
$p$ será verdadeira em todos os mundos distintos de~$w$ e, portanto, em todos
os mundos acessíveis a partir de~$w$ (pois está garantido que $w$ não é
acessível a partir de~$w$, e $w$ é o único mundo no qual~$p$ é falsa). Por
outro lado, $p$ é falsa em $w$, de modo que $\Box p \lif p$ é falsa em~$w$.

Isso sugere que devemos introduzir uma notação para estruturas de modelos sem
uma valoração: nós as chamamos de \emph{frames}. Um frame $\mModel{F}$ é
simplesmente um par $\tuple{W, R}$ formado por um conjunto de mundos com uma
relação de acessibilidade. Dizemos, então, que todo modelo $\tuple{W, R, V}$
é \emph{baseado no} frame $\tuple{W, R}$. Reciprocamente, um frame determina
a classe de modelos nele baseados; e uma classe de frames determina a classe
de modelos baseados em algum frame da classe. Podemos definir
$\mModel{F} \Entails !A$, a noção de !!a{formula} ser \emph{válida} em um
frame, da seguinte maneira: $\mSat{M}{!A}$ para todo $\mModel{M}$ baseado
em~$\mModel{F}$.

Com essa notação, podemos estabelecer relações de correspondência entre
!!{formula}s e classes de frames: p.ex.,
$\mModel{F} \Entails \Box p \lif p$ se, e somente se, $\mModel{F}$ é
reflexivo.

\end{document}
